\documentclass[12pt]{article}

% --- Packages ---
\usepackage[margin=1in]{geometry}
\usepackage{setspace}
\usepackage{titlesec}
\usepackage{parskip}
\usepackage{hyperref}
\usepackage{listings}
\usepackage{xcolor}
\usepackage{booktabs}
\usepackage{caption}
\usepackage{fancyhdr}
\usepackage{amsmath}

% --- Hyperlink styling ---
\hypersetup{
    colorlinks=true,
    linkcolor=black,
    urlcolor=blue,
    citecolor=black
}

% --- Code listing style ---
\lstdefinestyle{gostyle}{
    language=Go,
    basicstyle=\ttfamily\footnotesize,
    keywordstyle=\color{blue}\bfseries,
    commentstyle=\color{gray}\itshape,
    stringstyle=\color{teal},
    numberstyle=\tiny\color{gray},
    numbers=left,
    stepnumber=1,
    frame=single,
    breaklines=true,
    captionpos=b,
    tabsize=4,
    showstringspaces=false,
    morekeywords={func, type, struct, var, const, package, import, return, for, range, if, else}
}

% --- Header / Footer ---
\pagestyle{fancy}
\fancyhf{}
\rhead{Progress Report 2}
\lhead{CSCI-411-01}
\cfoot{\thepage}

% --- Section formatting ---
\titleformat{\section}{\large\bfseries}{}{0em}{}[\titlerule]
\titleformat{\subsection}{\normalsize\bfseries}{}{0em}{}

% ============================================================
\begin{document}

% --- Title Page ---
\begin{titlepage}
    \centering
    \vspace*{2cm}

    {\LARGE\bfseries Progress Report 2\par}
    \vspace{0.5cm}
    {\large Empirical Comparison of Spatial Data Structures for\\
    Broad-Phase Collision Detection Under High Entity Churn\par}

    \vspace{2cm}

    \begin{tabular}{rl}
        \textbf{Author:}     & Sachin Pandit \\[6pt]
        \textbf{Course:}     & CSCI-411-01   \\[6pt]
        \textbf{Date:}       & September 28, 2026 \\[6pt]
        \textbf{Repository:} & \url{https://github.com/sachinpandit140/Collision-Evaluation-Study} \\
    \end{tabular}

    \vfill

    \textit{Department of Computer Science}

\end{titlepage}

% ============================================================
\setstretch{1.2}

% ------------------------------------------------------------
\section{Project Overview}

This project is an empirical comparative analysis of two spatial data structures used in broad-phase collision detection: \textbf{Bounding Volume Hierarchies (BVH)} with incremental refitting, and \textbf{K-Dimensional Trees (KD-Trees)} with spatial rebuilding. Both acceleration structures are benchmarked against an $O(N^2)$ brute-force baseline across high entity churn conditions ($1{,}000$ to $20{,}000$ continuously moving entities).

\textbf{Core Research Question:} \textit{What are the exact performance trade-offs between dynamic tree-rebuilding overhead and collision traversal efficiency under high entity churn?}

The primary metrics evaluated include per-frame tree refit/rebuild time, query traversal latency, $p_{99}$ latency spikes (frame stutter detection), and heap allocation pressure. All implementations are written in Go from scratch with zero third-party dependencies, enforcing strict execution determinism and cache-friendly contiguous memory layouts. The project repository is available at:
\begin{center}
    \url{https://github.com/sachinpandit140/Collision-Evaluation-Study}
\end{center}

% ------------------------------------------------------------
\section{Work Completed Since Report 1}

\subsection{Extended Literature Review}

Four additional publications were reviewed to prepare for the BVH and KD-Tree implementation phases:

\begin{enumerate}
    \item \textbf{Goldsmith \& Salmon (1987)} --- Proposed the original formulation of the Surface Area Heuristic (SAH) for ray-object intersection acceleration \cite{goldsmith1987}. Their cost model estimates traversal probability at each node by the ratio of child surface area to parent surface area. This directly informs the split-cost evaluator needed for top-down BVH construction in Phase~3: at each internal node, the build selects the axis and partition that minimizes the SAH cost $C = C_{\text{trav}} + \frac{SA(L)}{SA(P)} \cdot n_L \cdot C_{\text{isect}} + \frac{SA(R)}{SA(P)} \cdot n_R \cdot C_{\text{isect}}$.

    \item \textbf{Karras (2012)} --- Presented a parallel radix-tree algorithm for BVH construction on GPUs \cite{karras2012}. While this project targets single-threaded CPU execution, Karras's analysis of Morton-code linearization and its relationship to spatial locality is directly relevant: Morton codes map 2D coordinates into a 1D ordering that preserves spatial proximity, which can serve as a fast initial sort before SAH-based partitioning. The paper also establishes that tree quality (measured by SAH cost) and construction speed are fundamentally at tension, a trade-off this study aims to quantify empirically.

    \item \textbf{Lauterbach et al.\ (2009)} --- Evaluated LBVH (Linear BVH) construction, in which primitives are sorted by Morton code and split at bit boundaries rather than evaluated with a full SAH sweep \cite{lauterbach2009}. LBVH builds run in $O(N \log N)$ time (dominated by the sort) compared to $O(N^2)$ for exhaustive SAH evaluation. Their measurements show that LBVH tree quality is within 10--15\% of full SAH trees for uniformly distributed geometry, but degrades under high spatial clustering. This informs the experimental design: the benchmark must include both uniform and clustered entity distributions to expose quality differences between build strategies.

    \item \textbf{van den Bergen (1997)} --- Described an AABB tree with lazy incremental updates for deformable models \cite{vandenbergen1997}. Rather than rebuilding the full tree each frame, van den Bergen enlarges (``fattens'') node AABBs by a margin proportional to entity velocity. Refitting then only triggers when an entity exits its fattened bounds. This lazy-refit strategy reduces tree update cost from $O(N)$ to amortized $O(k)$ where $k$ is the number of entities whose fattened bounds were violated. This approach will be implemented as a configuration option for the BVH engine to measure the refit-frequency vs.\ tree-quality trade-off.
\end{enumerate}

\subsection{Implementation Progress}

Since Report~1, the following components were completed:

\begin{itemize}
    \item \textbf{Dynamic BVH engine (\texttt{pkg/bvh})}: Implemented a binary BVH with top-down SAH-based construction. At each internal node, the builder evaluates all split candidates across both axes by sweeping cumulative bounding boxes left-to-right and right-to-left, selecting the partition that minimizes $C = \frac{SA(L)}{SA(P)} \cdot n_L + \frac{SA(R)}{SA(P)} \cdot n_R$. Leaf nodes hold a single entity index. The tree uses a contiguous \texttt{[]node} slice (pre-allocated to $2N - 1$ capacity) for cache-friendly traversal. Bottom-up \texttt{Refit} updates bounding boxes without changing tree topology. Self-collision traversal via \texttt{FindCollisions} prunes subtree pairs whose bounding boxes do not overlap, descending into the larger node at each step.

    \item \textbf{KD-Tree engine (\texttt{pkg/kdtree})}: Implemented a spatial-median KD-Tree with alternating axis splits ($\text{depth} \bmod 2$). Entities are sorted by AABB center on the split axis and partitioned at the median index. The tree is fully rebuilt from scratch each frame, matching the standard KD-Tree usage pattern for dynamic scenes. The collision query uses the same self-collision/cross-collision traversal pattern as the BVH.

    \item \textbf{Multi-method benchmark harness (\texttt{cmd/bench/main.go})}: Extended the CLI runner with a \texttt{-method} flag accepting \texttt{brute}, \texttt{bvh}, or \texttt{kdtree}. The BVH path builds the tree on the first frame and refits on subsequent frames; the KD-Tree path rebuilds every frame; the brute-force path runs the $O(N^2)$ baseline. Tree construction/refit time is captured in the tree update phase; collision query time is measured separately.

    \item \textbf{Correctness verification}: All three methods produce identical collision pair counts on every frame under identical parameters (seed = 42, $N = 1{,}000$, 100 frames), confirming that both acceleration structures agree with the brute-force oracle. The test suite expanded from 18 to 32 unit tests across 5 packages, all passing.

    \item \textbf{Per-frame telemetry pipeline (\texttt{pkg/sim/metrics.go})}: Implemented a \texttt{Recorder} that accumulates per-frame \texttt{FrameRecord} structs capturing entity count, tree update time, query time, total frame time, heap allocation delta, GC pause time, and candidate pair count. The recorder computes $p_{50}$, $p_{95}$, and $p_{99}$ percentiles over total frame time and exports all frame data to structured CSV via \texttt{WriteCSV}.
\end{itemize}

% ------------------------------------------------------------
\section{Evidence of Progress}

\subsection{Source Code Implementations}

The following snippets illustrate the core algorithms implemented for the BVH and KD-Tree engines:

\begin{lstlisting}[style=gostyle, caption={SAH-based BVH construction with sweep evaluation (\texttt{pkg/bvh/bvh.go}, excerpt)}]
func (t *Tree) buildRecursive(entities []sim.Entity, indices []int32) int32 {
    if len(indices) == 1 {
        idx := int32(len(t.nodes))
        t.nodes = append(t.nodes, node{
            box: entities[indices[0]].Bounds,
            left: -1, right: -1, entityIdx: indices[0],
        })
        return idx
    }

    // Compute parent bounds
    parentBox := entities[indices[0]].Bounds
    for _, idx := range indices[1:] {
        parentBox = parentBox.Merge(entities[idx].Bounds)
    }
    parentArea := parentBox.Area()
    if parentArea < 1e-10 { parentArea = 1e-10 }

    bestCost := math.MaxFloat64
    bestAxis, bestSplit := -1, -1

    for axis := 0; axis < 2; axis++ {
        sortByAxis(entities, indices, axis)
        // Sweep cumulative AABBs left-to-right and right-to-left
        leftBoxes := make([]geom.AABB2, len(indices))
        curLeft := entities[indices[0]].Bounds
        leftBoxes[0] = curLeft
        for i := 1; i < len(indices); i++ {
            curLeft = curLeft.Merge(entities[indices[i]].Bounds)
            leftBoxes[i] = curLeft
        }
        rightBoxes := make([]geom.AABB2, len(indices))
        curRight := entities[indices[len(indices)-1]].Bounds
        rightBoxes[len(indices)-1] = curRight
        for i := len(indices) - 2; i >= 0; i-- {
            curRight = curRight.Merge(entities[indices[i]].Bounds)
            rightBoxes[i] = curRight
        }
        // Evaluate SAH cost at each split position
        for i := 0; i < len(indices)-1; i++ {
            nL, nR := float64(i+1), float64(len(indices)-i-1)
            cost := (leftBoxes[i].Area()/parentArea)*nL +
                    (rightBoxes[i+1].Area()/parentArea)*nR
            if cost < bestCost {
                bestCost, bestAxis, bestSplit = cost, axis, i+1
            }
        }
    }
    sortByAxis(entities, indices, bestAxis)
    // Reserve node, recurse into children
    nodeIdx := int32(len(t.nodes))
    t.nodes = append(t.nodes, node{box: parentBox, entityIdx: -1})
    left := t.buildRecursive(entities, indices[:bestSplit])
    right := t.buildRecursive(entities, indices[bestSplit:])
    t.nodes[nodeIdx].left = left
    t.nodes[nodeIdx].right = right
    return nodeIdx
}
\end{lstlisting}

\begin{lstlisting}[style=gostyle, caption={KD-Tree spatial-median construction (\texttt{pkg/kdtree/kdtree.go}, excerpt)}]
func (t *Tree) buildRecursive(entities []sim.Entity, indices []int32, depth int) int32 {
    if len(indices) == 1 {
        t.nodes = append(t.nodes, node{
            box: entities[indices[0]].Bounds,
            left: -1, right: -1, entityIdx: indices[0],
        })
        return int32(len(t.nodes) - 1)
    }

    axis := depth % 2 // alternate X, Y
    sort.Slice(indices, func(i, j int) bool {
        return center(entities[indices[i]].Bounds, axis) <
               center(entities[indices[j]].Bounds, axis)
    })
    mid := len(indices) / 2

    nodeIdx := int32(len(t.nodes))
    t.nodes = append(t.nodes, node{left: -1, right: -1, entityIdx: -1})

    leftIdx := t.buildRecursive(entities, indices[:mid], depth+1)
    rightIdx := t.buildRecursive(entities, indices[mid:], depth+1)

    box := t.nodes[leftIdx].box.Merge(t.nodes[rightIdx].box)
    t.nodes[nodeIdx].box = box
    t.nodes[nodeIdx].left = leftIdx
    t.nodes[nodeIdx].right = rightIdx
    return nodeIdx
}
\end{lstlisting}

\subsection{Benchmark Telemetry}

All three collision detection methods were executed on an Apple M4 system with fixed parameters ($\text{seed}=42$, world dimensions $1000 \times 1000$, entity size $2 \times 2$, 100 frames). Table~\ref{tab:benchmarks} presents the comparative measurements at $N = 1{,}000$:

\begin{table}[ht]
\centering
\small
\begin{tabular}{lcccc}
\toprule
\textbf{Method} & \textbf{$p_{50}$ (ms)} & \textbf{$p_{95}$ (ms)} & \textbf{$p_{99}$ (ms)} & \textbf{Speedup vs.\ Brute} \\
\midrule
Brute-force $O(N^2)$ & 1.874 & 3.116 & 4.293 & $1.0\times$ \\
BVH (SAH + refit)     & 0.059 & 0.099 & 0.105 & $31.8\times$ \\
KD-Tree (rebuild)     & 0.458 & 0.609 & 0.851 & $4.1\times$ \\
\bottomrule
\end{tabular}
\caption{Comparative performance at $N = 1{,}000$ over 100 frames. All three methods produce identical collision pairs on every frame (751 total pairs across 100 frames), confirming correctness. Speedup is computed from $p_{50}$ ratios.}
\label{tab:benchmarks}
\end{table}

At $N = 1{,}000$, the BVH with incremental refitting is $31.8\times$ faster than brute force at the median, reflecting the combined benefit of SAH-optimized tree quality and $O(N)$ refit cost on subsequent frames. The KD-Tree, which pays the full $O(N \log N)$ rebuild cost every frame, still achieves a $4.1\times$ speedup over brute force but is $7.8\times$ slower than the BVH. This gap directly validates the central research hypothesis: at moderate entity counts, refit-based maintenance dominates rebuild-based maintenance.

The $p_{99}$ tail latency gap is also notable: the BVH's $p_{99}$ of 0.105~ms is only $1.8\times$ its $p_{50}$, indicating stable frame-to-frame performance. The brute-force $p_{99}$ of 4.293~ms is $2.3\times$ its $p_{50}$, reflecting higher variance from cache pressure during the $O(N^2)$ sweep.

All 32 automated unit tests pass with zero failures:
\begin{center}
\texttt{PASS: pkg/geom (8) \quad PASS: pkg/sim (5) \quad PASS: pkg/baseline (5) \quad PASS: pkg/bvh (7) \quad PASS: pkg/kdtree (7)}
\end{center}

% ------------------------------------------------------------
\section{Challenges}

\textbf{BVH refit quality degradation:} The current BVH refits bounding boxes without restructuring the tree topology. Over many frames, entity movement causes internal node bounding boxes to expand monotonically, reducing the pruning power of overlap checks during traversal. At $N = 1{,}000$ over 100 frames this effect is not yet visible in the timing data, but longer simulations and higher entity counts will expose it. The van den Bergen (1997) fattened-AABB approach is a planned mitigation: enlarging node bounds by a velocity-proportional margin and only triggering refit when an entity exits its fattened region.

\textbf{KD-Tree rebuild cost dominance:} The KD-Tree's $O(N \log N)$ full rebuild accounts for the majority of its per-frame cost, as reflected in the tree update phase timing. At $N = 1{,}000$, the sort-dominated rebuild is still fast enough to beat brute force, but the gap narrows at higher entity counts. Partial-rebuild strategies (reconstructing only subtrees whose bounding volume expanded beyond a threshold) are worth evaluating, though they add implementation complexity.

\textbf{Measurement noise from \texttt{runtime.ReadMemStats}:} Calling \texttt{ReadMemStats} triggers a stop-the-world pause in Go to collect accurate heap statistics. This injects $\sim$0.1~ms of overhead per frame into the timing data. The final benchmark suite will use a post-hoc sampling strategy (reading stats every $k$-th frame) to reduce measurement interference.

% ------------------------------------------------------------
\section{Next Steps}

With Phases 1--4 complete (primitives, simulation engine, brute-force baseline, BVH, and KD-Tree), the remaining milestones are:

\begin{enumerate}
    \item \textbf{Multi-$N$ benchmark matrix}:
    Sweep entity counts from $1{,}000$ to $20{,}000$ in increments of $1{,}000$, running all three methods under identical seeds. The target is to identify the crossover point where BVH refit quality degrades below KD-Tree rebuild quality, and to characterize the scaling exponents of each method empirically.

    \item \textbf{Phase 5 --- WASM/Canvas visualizer (\texttt{cmd/visualizer}, \texttt{web/})}:
    Compile the simulation engine to WebAssembly and render entity AABBs, BVH node wireframes (color-coded by tree depth), and KD-Tree split planes on an HTML5 Canvas. Add live metric readouts (frame time, candidate pairs, tree depth) as a debug overlay.

    \item \textbf{Final benchmark suite and paper data collection}:
    Run the full comparative benchmark matrix, export structured CSV telemetry, and generate publication-quality plots (matplotlib/seaborn) for the final paper submission.
\end{enumerate}

% ------------------------------------------------------------
\begin{thebibliography}{9}

\bibitem{bentley1975}
Bentley, J.~L. (1975).
\newblock Multidimensional binary search trees used for associative searching.
\newblock \textit{Communications of the ACM}, 18(9), 509--517.
\newblock \url{https://doi.org/10.1145/361002.361007}

\bibitem{wald2007}
Wald, I. (2007).
\newblock On fast construction of SAH-based bounding volume hierarchies.
\newblock In \textit{Proceedings of the 2007 IEEE Symposium on Interactive Ray
Tracing (RT '07)}.
\newblock \url{https://doi.org/10.1109/RT.2007.4342588}

\bibitem{kopta2012}
Kopta, D., Ize, T., Spjut, J., Brunvand, E., Davis, A., \& Kensler, A. (2012).
\newblock Fast, effective BVH updates for animated scenes.
\newblock In \textit{Proceedings of the ACM SIGGRAPH Symposium on Interactive
3D Graphics and Games (I3D '12)}.
\newblock \url{https://doi.org/10.1145/2159616.2159632}

\bibitem{ericson2005}
Ericson, C. (2005).
\newblock \textit{Real-Time Collision Detection}.
\newblock Morgan Kaufmann.
\newblock ISBN: 978-1558607323.

\bibitem{goldsmith1987}
Goldsmith, J. \& Salmon, J. (1987).
\newblock Automatic creation of object hierarchies for ray tracing.
\newblock \textit{IEEE Computer Graphics and Applications}, 7(5), 14--20.
\newblock \url{https://doi.org/10.1109/MCG.1987.276983}

\bibitem{karras2012}
Karras, T. (2012).
\newblock Maximizing parallelism in the construction of BVHs, octrees, and $k$-d trees.
\newblock In \textit{Proceedings of the Fourth ACM SIGGRAPH / Eurographics Conference on High-Performance Graphics (HPG '12)}.
\newblock \url{https://doi.org/10.2312/EGGH/HPG12/033-037}

\bibitem{lauterbach2009}
Lauterbach, C., Garland, M., Sengupta, S., Luebke, D., \& Manocha, D. (2009).
\newblock Fast BVH construction on GPUs.
\newblock \textit{Computer Graphics Forum}, 28(2), 375--384.
\newblock \url{https://doi.org/10.1111/j.1467-8659.2009.01377.x}

\bibitem{vandenbergen1997}
van den Bergen, G. (1997).
\newblock Efficient collision detection of complex deformable models using AABB trees.
\newblock \textit{Journal of Graphics Tools}, 2(4), 1--13.
\newblock \url{https://doi.org/10.1080/10867651.1997.10487480}

\end{thebibliography}

\end{document}
